\begin{table*}[t]
\centering
\small
\begin{tabular}{@{}lcccccccc@{}}
\toprule
& \multirow{2}{*}{\textbf{Target}} & \multicolumn{2}{c}{\textbf{Multi-Constraint (tpwgts)}} & \multicolumn{2}{c}{\textbf{Multi-Constraint (ubvec)}} & \multicolumn{3}{c}{\textbf{Edge-Weighted (Optimal)}} \\
\cmidrule(lr){3-4} \cmidrule(lr){5-6} \cmidrule(lr){7-9}
\textbf{State (Minority \%)} & \textbf{MM} & MM & Max \% & MM & Max \% & Config & MM & Max \% \\
\midrule
Mississippi (46.1\%) & 2 & \textbf{2} & 53.3 & \textbf{2} & 53.2 & 1x @ 40\% & \textbf{2} & 53.8 \\
Georgia (42.4\%) & 5 & \textbf{5} & 70.4 & \textbf{5} & 76.7 & 1x @ 40\% & \textbf{5} & 76.9 \\
Louisiana (41.6\%) & 2 & 1 & 58.8 & 1 & 52.9 & 5x @ 40\% & \textbf{2} & 55.8 \\
Alabama (36.9\%) & 2 & 0 & 47.3 & 0 & 49.6 & 5x @ 40\% & \textbf{2} & 50.8 \\
South Carolina (35.1\%) & 3 & 0 & 44.4 & 0 & 47.2 & 50x @ 50\% & 0 & 49.2 \\
\midrule
\textbf{Success Rate} & & \multicolumn{2}{c}{40\% (2/5)} & \multicolumn{2}{c}{40\% (2/5)} & & \multicolumn{2}{c}{\textbf{80\% (4/5)}} \\
\bottomrule
\end{tabular}
\caption{Comprehensive comparison of VRA compliance approaches across five covered states. Bold indicates meeting or exceeding target MM district count. Multi-constraint methods achieve 40\% success rate, while edge-weighted optimization doubles this to 80\%. Alabama represents the critical breakthrough: 0 MM districts with multi-constraint vs 2 MM districts with edge-weighting. Optimal edge-weighted configuration shown (lowest weight factor achieving success).}
\label{tab:comprehensive_comparison}
\end{table*}
